% =====================================================================
%  Section 13 — Phase 2 questions conclusion
% =====================================================================

\section{Phase 2 \textbullet\ What the Thirteen Questions Collectively Show}
\label{sec:phase2-questions-conclusion}

The thirteen Phase 2 research questions are not independent — they are
a single argument disguised as a sequence. Read together, they
establish four claims that the formal hypothesis tests in
Section~\ref{sec:phase2-hypotheses} will then formalise.

\subsection*{Claim 1 — The new infrastructure does not follow density}

Q13 (per-line scatter), Q14 (84\% of ghosts beyond 2\,km of any new
metro), Q16 (fastest-growing districts with zero coverage), and Q22
(top-15 worst residuals all in Giza/Qalyubia) all point in the same
direction: the post-2014 metro/LRT investment lands away from where
people already live.

\subsection*{Claim 2 — Underservedness is structural, not just dense}

Q17 ($\rho = 0.548$) shows that the densest hexes are systematically
the most underserved. Q18 ($\rho = 0.025$, flat) shows that informal
modal share is the same across density tiers, which means the
informal network is the substrate everywhere, not a poverty pocket.
Together: Cairo has an everywhere-informal background and a
selectively-formal foreground that does not align.

\subsection*{Claim 3 — BRT is the contrast case}

Q20 (BRT corridor demand) shows that BRT stations sit on top of
\emph{existing} informal demand. Q18b's gap-closure matrix bears it
out: the only cell at 25\% (the matrix maximum) is BRT $\times$ G3
(vehicle mismatch). \kw{Not all post-2014 infrastructure is equally
weak — BRT is what demand-aligned planning looks like in this dataset.}

\subsection*{Claim 4 — There is a real, sized market for closing the gap}

Q24's K-Means delivers a 4-cluster partition with \stat{ARI = 1.00}
stability. Three of those clusters
(Formal-Served Core + Peripheral Growth + Mixed) make up
\stat{62 districts, 18.9 million residents} —
the addressable market for a product that visibly bridges the two
systems. Q25 shows the bridges already exist physically; what is
missing is the software layer that surfaces them.

\subsection*{The market-sizing visual}

\figitem{phase2/market_sizing_bar.png}{1.0}{Cluster population scale —
Formal-Served Core, Peripheral Growth, and Mixed sum to 18.9M
residents. The Low-Activity Outskirts cluster is excluded from the
addressable market.}{fig:market-sizing}

\figitem{phase2/sunburst_market.png}{1.0}{Sunburst — governorate $\to$ K-Means cluster $\to$ district, sized by 2023 population. Reveals the geographic distribution of each market segment.}{fig:sunburst}

\subsection*{What the questions cannot do alone}

Descriptive questions describe; they do not formally test. The
quantitative claims in this section are visible in the Plotly figures
but they remain \emph{visual} arguments. The next section
(\ref{sec:phase2-hypotheses}) takes the three strongest of these
visual arguments — coverage--need mismatch, LRT catchment deficit, BRT
corridor match — and runs each of them through a non-parametric
hypothesis test with effect sizes. That is where ``the dense ones
look underserved'' becomes \stat{$\varepsilon^2 = 0.826$} and ``LRT
serves few people'' becomes \stat{Cliff's $\delta = -0.991$}.
